% !TEX root = ../../../tree_notes.tex
% =====================================================================
%  13.1  Tree DP: state per node (subtree DP)  sec:c10-subtrees
%  Exports: def:c10-subtree-state  thm:c10-subtree-template
%           lst:c10-fold  fig:c10-tuples  ex:c10-subtrees-run
%           prob:c10-subtree-solved (LC 333)  lst:c10-subtree-soln
%  The general-tree sibling of Ch.1's thm:c01-height-recursion and the
%  tuple-carrying upgrade of Ch.12's lst:c09-dfs3: same DFS skeleton,
%  richer return value.  Everything below is hand-computed on T.
% =====================================================================
\section{Subtree DP: One Tuple per Node, Folded from Children}\label{sec:c10-subtrees}

\begin{previouslybox}
\begin{enumerate}[nosep, leftmargin=1.6em]
  \item Chapter~1 proved $\height(v)=1+\max_c\height(c)$ and
        $\size(v)=1+\sum_c\size(c)$ by induction on the height of $v$'s
        subtree (\thmref{c01-height-recursion}). Which two conventions made
        a leaf fall out of the same formula without an \code{if}?
        ($\height(\nil)=-1$, $\size(\nil)=0$)
  \item Chapter~2's bottom-up pattern: what does a node ``read'' from its
        children, and in what order does the read happen?
        (return values, after every child has returned; the fold is the
        unwind)
  \item Chapter~12's \coderef{c09-dfs3} pushed $\depth$ \emph{down} as an
        argument and folded $\height$ and $\size$ \emph{up}. What
        \emph{three} properties did that recursion maintain, and does the
        same skeleton survive if we replace the fold with an arbitrary
        tuple? (yes -- the parent guard and the post-order are the only
        dependencies)
  \item Chapter~4 rooted LC~2049's ``count nodes by same-label subtree''
        problem via a parent parameter and one DFS. When does \emph{that}
        pattern stop being enough and start needing a \emph{tuple} at every
        node? (as soon as one number cannot carry the answer -- the answer
        needs at least two coordinates to close under the fold)
\end{enumerate}
\end{previouslybox}

\begin{topicintro}
Chapter~12's \coderef{c09-dfs3} was already a DP: each node folded its
children's numbers into two scalars. Chapter~1's \thmref{c01-height-recursion}
was the same idea on binary trees. Both use the \emph{same} DFS skeleton --
parent guard, post-order, one number per node -- and differ only in
\emph{what} they carry. This section makes that pattern explicit and
general: the state at a node is a \emph{tuple} $\mathcal S(v)$ chosen so that
$\mathcal S(v)$ is a computable function of $\{\mathcal S(c):c\in\children(v)\}$
alone. Once the tuple is picked, the algorithm is the same three lines
everywhere. The design question -- \emph{what tuple closes the recurrence} --
is where all the difficulty lives, and the LC~333 subtree-BST tuple
(\probref{c10-subtree-solved}) is a first exercise in picking one.

\textbf{Prerequisites.} \thmref{c01-height-recursion} (Ch.1) and
\coderef{c09-dfs3} (Ch.12) -- we reuse both verbatim, we do not re-derive
either; \chref{c04} for rooting by parent parameter; the bottom-up habit of
\chref{c02}.

\textbf{After this section you can}
\begin{itemize}[nosep]
  \item state the child-fold template $\mathcal S(v)=\mathcal F(v,\,\text{children})$
        and rewrite \coderef{c09-dfs3} in it; \hfill\lv{2}
  \item pick a tuple for a given problem statement and argue why it is
        the smallest tuple that closes; \hfill\lv{3}
  \item cold-write LC~333's subtree-BST tuple and its trace; \hfill\lv{3}--\lv{4}
\end{itemize}
\end{topicintro}

\subsection{Intuition}

A binary tree stores ``left'' and ``right'' inside the data structure; a
general tree stores only ``neighbours'' (\chref{c04}). Once we root by parent
parameter, ``children'' becomes a set $\children(v)$ of arbitrary size, and
every recurrence that used to read two named slots must be rewritten to fold
a \emph{set}. Chapter~1's $\max(h_L,h_R)+1$ becomes $1 + \max_{c}\height(c)$;
Chapter~1's $\size_L+\size_R+1$ becomes $1 + \sum_c\size(c)$; the initialisers
$\height(\nil)=-1$ and $\size(\nil)=0$ turn into the $\max$ and $\sum$ identity
elements ($0$ for height when there are no children, $1$ for size once we
count $v$ itself). Nothing new is proved -- the induction is the same
$\mathit{dfs}(u,p)$/children-first argument from \thmref{c09-onepass};
only the \emph{payload} changes.

The payload upgrade is what turns the metric chapter into a DP chapter.
A \emph{scalar} per node (Ch.12: depth, height, size) is enough for shape
measures because the answer is a $\min$, $\max$ or $\sum$ of scalars. The
moment the answer depends on \emph{multiple} facts about a subtree -- ``is it
a BST, and its min, and its max, and its size'' -- a scalar is not enough and
the tuple is where you write down what you must remember. Ch.2 was already
doing this on binary trees (House Robber III's pair (rob, skip) is a
two-coordinate tuple); this chapter is the general-tree version of the same
habit.

\subsection{Definitions and Notation}

\begin{definition}{Subtree state and child-fold}{c10-subtree-state}
Fix a root $r$ in a tree $T$ and let $\children(v)$ be the child set of $v$
in that rooted view (\defref{c09-measures}). A \term{subtree state} is a
function $\mathcal S\colon V\to\mathfrak T$ into some set $\mathfrak T$
defined by
\[
  \mathcal S(v) \;=\; \mathcal F\!\Big(v,\;\big\{\mathcal S(c):c\in\children(v)\big\}\Big),
\]
where $\mathcal F$ is a fixed \emph{fold rule} that depends only on $v$ and
the multiset of its children's states. The \term{answer} is a function
$A(\mathcal S(r))$ read at the root (or, for per-node outputs, $A(\mathcal S(v))$
for every $v$). We always require $\mathcal F$ to be well-defined on the empty
child set (a leaf's state is $\mathcal F(v,\emptyset)$, the tuple's
initialiser); this is the general-tree restatement of Ch.1's
$\height(\nil)=-1$ convention.
\end{definition}

\noindent\textbf{In plain words.} Each node carries \emph{one} tuple; the
tuple is computed by looking only at the tuples of its children (and the
node's own label, if any). Whatever the tuple is, the algorithm to fill it
everywhere is one DFS. \emph{Choosing} the tuple is the whole problem.

\begin{theorem}{The child-fold template computes $\mathcal S$ in $\bigO{n}$}{c10-subtree-template}
For any tree $T$ with $n$ nodes and any state $\mathcal S$ and fold $\mathcal F$
of \defref{c10-subtree-state} whose per-node cost is $\bigO{\deg(v)}$,
Algorithm \coderef{c10-fold} fills $\mathcal S(v)$ for every $v$ in
$\Theta(n)$ time and $\bigO{h}$ auxiliary stack space, and its output
satisfies the recurrence.
\end{theorem}
\begin{proofbox}
Structurally identical to \thmref{c09-onepass}. Each node is entered once
(unique path from $r$, parent guard rejects the return edge) and its loop
iterates over $\children(v)$, contributing $\deg(v)-1$ fold steps. Total
loop work is $\sum_v (\deg(v)-1) = 2(n-1) - (n-1) = n-1$ calls to the
per-child update -- $\Theta(n)$ overall. Correctness is strong induction on
the height of $v$'s subtree: $\mathit{dfs}(v,p)$ returns only after every
child call has returned, so by induction the multiset
$\{\mathcal S(c)\}$ is already correct, and $\mathcal F$ produces
$\mathcal S(v)$ by definition. Stack depth $=h$.
\end{proofbox}

\begin{insightbox}
The \emph{recurrence} is the entire problem. Given $\mathcal F$ and the
post-order DFS of \coderef{c09-dfs3}, the algorithm writes itself -- three
lines and one tuple. Designing $\mathcal F$ is the design of the DP.
Two failure modes to keep in mind: a tuple \emph{too small} silently breaks
the closure ($\mathcal F$ looks at coordinates it does not have); a tuple
\emph{too large} is correct but slow (worse $\bigO{\mathfrak T}$, sometimes
worse time). Ch.14 (\chref{c27}) makes the size-vs-cost trade precise for
budgeted merges; this chapter stays in the constant-tuple regime.
\end{insightbox}

\subsection{The Algorithm: The Child-Fold Template}

\begin{lstlisting}[caption={The generic child-fold on a general tree: same skeleton as \coderef{c09-dfs3}, arbitrary tuple payload}, label={lst:c10-fold}]
// adj[u] = neighbours; root at r; call fold_dfs(r, -1).
// State S[u] is any struct with an initialiser for a leaf.
vector<vector<int>> adj;
vector<State> S;                          // sized n

State fold_node(int v, const vector<State>& kids) {
    State s;                              // v's own label may seed s
    s = F(v, kids);                       // THE ONE FUNCTION TO DESIGN
    return s;
}

void fold_dfs(int v, int p) {
    vector<State> kids;
    for (int u : adj[v]) if (u != p) {    // parent guard = rooting
        fold_dfs(u, v);                   // child state final first...
        kids.push_back(S[u]);             // ...then collected
    }
    S[v] = fold_node(v, kids);            // fold closes the tuple
}
// after the call: answer is A(S[r]), or A(S[v]) read per node
\end{lstlisting}

\begin{complexitybox}
Time $\Theta(n)$ by \thmref{c10-subtree-template}, assuming $|F(v,\cdot)|$
is $\bigO{\deg(v)}$. Memory: $\bigO{n}$ for the state array plus $\bigO{h}$
for the call stack. Verdict: the whole chapter reduces to designing $F$ and
$A$ -- the skeleton is not the difficulty. When a chain is a legal input at
$n \ge 2\cdot 10^{5}$, substitute \coderef{c09-iter}'s two-phase form
verbatim; reversed preorder is still a safe fold order (Ch.12's justification
carries across unchanged).
\end{complexitybox}

\subsection{Worked Example: A Two-Coordinate Tuple on $T$}

\begin{example}{Subtree sizes and sum-of-descendant-depths on the running tree $T$}{c10-subtrees-run}
Take $T$ from \figref{c10-runtree-dp}, root at~$1$, and set
\[
  \mathcal S(v) \;=\; \big(\mathit{sz}(v),\;\mathit{LD}(v)\big),\quad
  \mathit{LD}(v) \;=\; \sum_{u\in\subtree(v)} d(v,u),
\]
the pair (subtree size, sum of depths \emph{measured from $v$}). The fold is
\[
  \mathcal S(v) \;=\; \Big(1+\!\!\sum_{c}\!\mathit{sz}(c),\;\;\sum_{c}\!\big(\mathit{LD}(c)+\mathit{sz}(c)\big)\Big),
\]
because every descendant of $c$ sits one edge deeper from $v$ than from $c$:
$\mathit{LD}$ picks up one per node in $c$'s subtree, hence $\mathit{LD}(c)+\mathit{sz}(c)$.
Table below is hand-computed on $T$ and reused \emph{unchanged} in
\secref{c10-twostate}, \secref{c10-harvest} and \secref{c10-distance}.
\end{example}

\begin{figure}[H]
\centering
\begin{tikzpicture}[diagram, scale=0.88, transform shape,
    tnode/.style={circle, draw=primaryblue, line width=0.9pt, fill=primaryblue!8,
                  minimum size=6.2mm, font=\small\bfseries, text=primaryblue!60!black}]
  \node[tnode] (n1) at (0,0) {$1$};
  \node[tnode] (n2) at (-3,-1.6) {$2$};
  \node[tnode] (n3) at ( 3,-1.6) {$3$};
  \node[tnode] (n4) at (-5.2,-3.2) {$4$};
  \node[tnode] (n5) at (-1.6,-3.2) {$5$};
  \node[tnode] (n6) at ( 1.6,-3.2) {$6$};
  \node[tnode] (n7) at ( 5.2,-3.2) {$7$};
  \node[tnode] (n8) at (-1.6,-4.8) {$8$};
  \node[tnode] (n9) at ( 5.2,-4.8) {$9$};
  \node[tnode] (n10) at ( 5.2,-6.4) {$10$};
  \draw[tedge] (n1)--(n2) (n1)--(n3) (n2)--(n4) (n2)--(n5)
               (n3)--(n6) (n3)--(n7) (n5)--(n8) (n7)--(n9) (n9)--(n10);
  \foreach \nd/\tr in {n1/10:20, n2/4:4, n3/5:7, n4/1:0, n5/2:1,
                      n6/1:0, n7/3:3, n8/1:0, n9/2:1, n10/1:0} {
    \node[font=\tiny\ttfamily, text=warningred!75!black, anchor=north]
      at (\nd.south) {(\tr)};}
  \node[annot, font=\scriptsize, anchor=south west, text=successgreen!50!black]
    at (-0.4,0.55) {read (sz : LD) below each node};
\end{tikzpicture}
\caption{Every node of $T$ after one pass of \coderef{c10-fold}, tagged
$(\mathit{sz},\mathit{LD})$ in red. Sanity: $\mathit{LD}(1)=20$ matches
Ch.12's hand-verified ``sum of depths at root $1$'' $=0\!+\!1\!+\!1\!+\!2\!+\!2\!+\!2\!+\!2\!+\!3\!+\!3\!+\!4=20$.}
\label{fig:c10-tuples}
\end{figure}

\begin{center}
\small
\renewcommand{\arraystretch}{1.15}
\begin{tabular}{l|cccccccccc}
\toprule
node $v$           & 1 & 2 & 3 & 4 & 5 & 6 & 7 & 8 & 9 & 10 \\
\midrule
$\mathit{sz}(v)$   & 10& 4 & 5 & 1 & 2 & 1 & 3 & 1 & 2 & 1 \\
$\mathit{LD}(v)$   & 20& 4 & 7 & 0 & 1 & 0 & 3 & 0 & 1 & 0 \\
\bottomrule
\end{tabular}
\end{center}

\begin{tracebox}[Fold order through \coderef{c10-fold} on $T$, root $1$, tuple $(\mathit{sz},\mathit{LD})$]
post-order (one legal): 4, 8, 5, 2, 10, 9, 7, 6, 3, 1.\\
4: children=() -> sz=1, LD=0. \;\; 8: children=() -> sz=1, LD=0.\\
5: children=(8) -> sz=1+1=2, LD=(0+1)=1.\\
2: children=(4,5) -> sz=1+1+2=4, LD=(0+1)+(1+2)=4.\\
10: children=() -> sz=1, LD=0. \;\; 9: children=(10) -> sz=2, LD=0+1=1.\\
7: children=(9) -> sz=3, LD=1+2=3. \;\; 6: children=() -> sz=1, LD=0.\\
3: children=(6,7) -> sz=1+1+3=5, LD=(0+1)+(3+3)=7.\\
1: children=(2,3) -> sz=1+4+5=10, LD=(4+4)+(7+5)=20 \checkmark (=n sanity,
LD(1)=sum of depths at root=20 from Ch.12's table). \checkmark
\end{tracebox}

\begin{insightbox}[Key takeaways]
\begin{enumerate}[nosep, leftmargin=1.6em]
  \item The skeleton is \coderef{c09-dfs3} again; only the fold changed.
        Every section of this chapter re-uses \lstinline{fold_dfs} verbatim.
  \item A two-coordinate tuple is \emph{cheaper than an extra pass}: computing
        $\mathit{LD}$ alone would need Ch.12's DFS plus a second accumulation,
        but folded together with $\mathit{sz}$ it costs one addition per edge.
  \item Sanity checks come from \emph{closed forms} you already know:
        $\mathit{LD}(r)=\sum_v\depth(v)$ (Ch.12) and $\mathit{sz}(r)=n$
        (Ch.1's $\size(r)=n$ convention) both fire on $T$.
\end{enumerate}
\end{insightbox}

\subsection{Variations and Twists}

\begin{itemize}[nosep]
  \item \textbf{Per-node answer, not root answer}: many problems ask
        ``for every $v$, report $\mathcal S(v)$''. Nothing changes in
        \coderef{c10-fold}; you just read $A(\mathcal S(v))$ after the DFS
        instead of only at $r$. LC~2049 (\chref{c04}) is exactly this shape
        with $\mathcal S=(\mathit{sz},\mathit{prod})$. \hfill\lv{3}
  \item \textbf{Tuple with a ``bad'' flag}: the LC~333 tuple below carries a
        boolean $\mathit{isBST}$ so a violation propagates up without
        requiring a re-visit; the same trick shows up in Ch.2's LC~814
        (prune-tree) as the ``keep-me'' bit. \hfill\lv{3}
  \item \textbf{Tuple size blows up}: when the fold must carry an
        \emph{array} (an $\bigO{k}$ histogram in \secref{c10-distance}, or a
        $\bigO{V}$ knapsack in Ch.14), the skeleton survives; the analysis
        does not. That is Chapter~14's (\chref{c27}) territory -- this chapter
        previews without proving. \hfill\lv{4}
  \item \textbf{Reroot preview}: the table above is \emph{for root $1$}. Asking
        $\mathcal S$ for every root is the rerooting question, Chapter~15
        (\chref{c20}) -- Ch.12's \secref{c09-radius} already previewed the
        distance-sum version and this chapter does not re-preview.
        \hfill\lv{4}
\end{itemize}

\subsection{Common Pitfalls}

\begin{warningbox}[Trap 1 -- tuple too small, fold fails silently]
The classic LC~333 mistake: carry only $(\mathit{isBST},\mathit{size})$ and
forget $(\mathit{minVal},\mathit{maxVal})$. A parent then tests ``my value is
bigger than everything on the left'' without a way to see what ``everything''
is -- the tuple closed under $\max$/$\sum$ folds in Ch.1 but does \emph{not}
close under BST-comparison folds. Symptom: your code returns the size of a
\emph{maximal} BST-shaped subtree but accepts a $\mathit{yes}$ at a node
whose left subtree is technically out-of-order at a deep leaf. Fix: extend
the tuple.
\end{warningbox}

\begin{warningbox}[Trap 2 -- reading the answer at every leaf, not the root]
When the problem asks for a global optimum (LC~333's ``largest BST
\emph{anywhere}'', Ch.12's tree diameter), reading $A(\mathcal S(r))$ misses
subtrees that are optimal at an internal node. The fix is \emph{not} to
re-run the DFS: extend $\mathcal S$ to carry ``best seen in my subtree'' as
an extra coordinate so the fold is
\[
  \mathit{best}(v) \;=\; \max\!\big(\text{is-BST-at-}v \;?\; \mathit{sz}(v):0,\;\;\max_c \mathit{best}(c)\big),
\]
and the final answer \emph{is} readable at $r$. This is the ``state + best''
pattern; it will reappear in \secref{c10-harvest} as the peak-sum and in
\secref{c10-distance} as the histogram-of-counts.
\end{warningbox}

\begin{warningbox}[Trap 3 -- trusting the recursion budget]
\thmref{c10-subtree-template}'s $\bigO{h}$ stack term is Ch.12's same trap
wearing a new hat: at $n\ge 2\cdot 10^{5}$ with chains allowed,
\coderef{c10-fold} overflows the call stack, and it is a \emph{runtime
error}, not a wrong answer -- invisible on hand samples. Substitute
\coderef{c09-iter}'s two-phase explicit-stack skeleton verbatim (Ch.12's
justification for reversed preorder carries across: in a tree, every child
is pushed only by its parent, so reversed preorder is a safe fold order).
\end{warningbox}

\subsection{Problem Lab}

The section's anchor is LeetCode~333, \probref{c10-subtree-solved}: given a
binary tree, find the size of the largest subtree that is a BST. ``Largest''
is by \emph{number of nodes}; ``subtree'' means any node together with all
its descendants -- a \emph{rooted} subtree, not a path. This is the tuple
fold in its purest form: a single scalar cannot carry BST-ness across a
parent boundary, so you must design the tuple, then the fold, and the DFS is
\coderef{c10-fold} with those two decisions filled in.

% ---------------------------------------------------------------------------
\solvedproblem{c10-subtree-solved}{Largest BST Subtree}{LeetCode 333}
  {https://leetcode.com/problems/largest-bst-in-binary-tree/}{LC Medium L3, premium-era}
  {subtree state tuple, BST verification, min/max propagation, best-in-subtree}

\begin{problemstatement}[Largest BST Subtree (premium-era; statement and samples verified against the live problem)]
Given the \texttt{root} of a binary tree, return \emph{the number of nodes of
the largest subtree that is a Binary Search Tree}. A \textbf{subtree} here
means a node together with \emph{all its descendants}; a \textbf{BST} is a
tree whose inorder traversal is strictly ascending (equivalently: for every
node, every value in its left subtree is $<$ its value, every value in its
right subtree is $>$ its value, recursively). Constraints: $1\le n\le 10^{4}$,
$-10^{4}\le \texttt{Node.val}\le 10^{4}$. \textbf{Example~1}: tree
\texttt{[10,5,15,1,8,null,7]} $\Rightarrow$ output \texttt{3} (the subtree
rooted at node~$5$, containing $\{1,5,8\}$, is BST and has three nodes; the
subtree rooted at $10$ is \emph{not} BST because the subtree rooted at $15$
is not -- $15$'s right child is $7$, which violates $15<\text{any in right}$).
\textbf{Example~2}: tree \texttt{[4,2,7,2,3,null,3]} $\Rightarrow$ output
\texttt{2} (the subtree $\{2,3\}$ rooted at node~$2$ is BST; the whole tree
is not -- the left~$2$ under root~$4$'s left~$2$ is equal to it, violating
strictness). \textbf{Example~3}: tree \texttt{[1]} $\Rightarrow$ output
\texttt{1} (a single node is BST).
\end{problemstatement}

\begin{attemptgate}[Attempt it yourself first -- \lv{3}, give yourself 25 minutes]
The trap is trying to write a scalar recursion. Say out loud what four
coordinates must be in your tuple before the fold can close.
\peek{When a child subtree is a BST, what does its parent need to know to
      decide its own BST-ness?}
\nudge{Not just its size -- it also needs the child's min, the child's max,
      and whether the child \emph{itself} is a BST (four coordinates).}
\fullhint{Pass~1: at every node return a tuple $(\mathit{isBST},\mathit{lo},
\mathit{hi},\mathit{sz})$ where $\mathit{lo},\mathit{hi}$ are the min and max
value in $v$'s subtree (\emph{whether} $v$ is BST or not; on non-BST subtrees
they still propagate correctly so the ancestor test stays total). Pass~2:
combine at node $u$ as $\mathit{isBST}(u) = \mathit{isBST}(L)\wedge\mathit
{isBST}(R)\wedge (\mathit{sz}(L)=0 \vee \mathit{hi}(L)<u\!\to\!\mathit{val})
\wedge (\mathit{sz}(R)=0 \vee u\!\to\!\mathit{val}<\mathit{lo}(R))$. Maintain
a running $\mathit{best}=\max(\mathit{best},\mathit{sz}(u))$ whenever
$\mathit{isBST}(u)$; or, if you prefer no globals, extend the tuple with a
$\mathit{best}$ coordinate folded by $\max$. Answer is $\mathit{best}$ at
the root; $0$ on an empty tree.}
\end{attemptgate}

\solstep{Understanding the Problem}
Two things must reconcile. (i)~``Is BST'' is a \emph{global} property of a
subtree -- it is not determined at $v$ alone but requires the \emph{entire}
left subtree to hold values strictly $<v$ and the entire right subtree to
hold values strictly $>v$. (ii)~The answer wants the size of the
\emph{largest} subtree satisfying (i) \emph{somewhere} in the tree, not just
at the root. A naive top-down re-verification at every node is $\bigO{n^2}$
in a chain and $\bigO{n\cdot h}$ in general; the tuple fold does it in
$\bigO{n}$ because BST-ness \emph{composes}: if $L$ and $R$ are already
labelled BST/non-BST with correct $(\mathit{lo},\mathit{hi},\mathit{sz})$,
testing $v$ takes $\bigO{1}$.

\solstep{Key Observations}
\begin{observation}
BST-ness at $v$ is a function of $\big(\mathit{isBST}(L),\mathit{isBST}(R),
\mathit{hi}(L),\mathit{lo}(R),v\!\to\!\mathit{val}\big)$ only -- it does not
need to inspect any deeper value, so long as $L$ and $R$'s tuples correctly
carry their own min and max.
\end{observation}
\begin{observation}
Even on a \emph{non-BST} subtree, $\mathit{lo}$ and $\mathit{hi}$ are well
defined (they are just min and max over a set of values); propagating them
upward keeps the ancestor test total, at the cost of one extra $\max/\min$
per fold. The $\mathit{isBST}$ boolean is what actually blocks the answer.
\end{observation}
\begin{observation}
The problem wants a $\max$ over \emph{all} nodes of $\mathit{sz}$ filtered by
$\mathit{isBST}$. Extend the tuple with a fifth coordinate $\mathit{best}$
folded by $\max(\text{self-if-BST}, \mathit{best}(L), \mathit{best}(R))$ so
the answer is readable at $r$ and no global variable is needed.
\end{observation}

\solstep{Think It Through With Me}
\begin{thinkbox}
\deadend{Inorder traverse the whole tree once; the longest contiguous
ascending interval is the answer.}
{Inorder of a tree gives a linear sequence whose contiguous segments do
\emph{not} correspond to subtrees -- you would need ``contiguous AND
descendant-closed'', which is exactly re-stating the problem and costs an
extra $\bigO{n\cdot h}$ scan. It also breaks on a chain where the
descendant-closed sets are prefixes of the inorder.}
\deadend{At every node $v$, call a ``validate-BST'' that checks the whole
subtree rooted at $v$ in $\bigO{\text{subtree-size}}$ by walking it.}
{$\bigO{n^2}$ worst case (chain: root validate is $\bigO{n}$, its child is
$\bigO{n-1}$, ...; sum $=\Theta(n^2)$; $n=10^{4}$ makes it $10^{8}$ --
borderline TLE). Symptom: AC on samples, TLE on hidden. The tuple fixes it
by \emph{caching} the summary at every node so each node's check is $\bigO{1}$.}
\aha{Five-coordinate tuple $(\mathit{isBST},\mathit{lo},\mathit{hi},
\mathit{sz},\mathit{best})$, one post-order fold, $\bigO{n}$ total, answer
$=\mathit{best}$ at the root. Everything below is bookkeeping.}
\end{thinkbox}

\solstep{Solution}
Run \coderef{c10-fold} on the tree, storing per-node $\mathcal S(v)$ in a
five-field struct. At $v$ with children $L,R$ (either may be $\nil$, treated
as $(\mathit{isBST}{=}\text{true},\;\mathit{lo}{=}+\infty,\;\mathit{hi}{=}-\infty,\;\mathit{sz}{=}0,\;\mathit{best}{=}0)$,
Ch.1's convention $\height(\nil)=-1$ restated for this tuple):
$\mathit{lo}(v)=\min(\mathit{lo}(L),\,v\!\to\!\mathit{val},\,\mathit{lo}(R))$,
$\mathit{hi}(v)=\max(\mathit{hi}(L),\,v\!\to\!\mathit{val},\,\mathit{hi}(R))$,
$\mathit{sz}(v)=1+\mathit{sz}(L)+\mathit{sz}(R)$,
$\mathit{isBST}(v)=\mathit{isBST}(L)\wedge\mathit{isBST}(R)\wedge\mathit{hi}(L)<v\!\to\!\mathit{val}<\mathit{lo}(R)$,
and $\mathit{best}(v)=\max\big(\mathit{isBST}(v)?\mathit{sz}(v):0,\;\mathit{best}(L),\;\mathit{best}(R)\big)$.
The answer is $\mathit{best}(r)$.

\solstep{Proof of Correctness}
\begin{proofbox}[The tuple returned at $v$ satisfies \defref{c10-subtree-state}, and $\mathit{best}(r)$ is the size of the largest BST subtree]
By strong induction on the height of $v$'s subtree. \emph{Base}: a leaf has
$\mathit{sz}{=}1$, $\mathit{lo}{=}\mathit{hi}{=}v\!\to\!\mathit{val}$, and
$\mathit{isBST}{=}\text{true}$ (the two boundary inequalities hold because
$\nil$'s $\mathit{hi}=-\infty$ and $\mathit{lo}=+\infty$); $\mathit{best}{=}1$
is the correct answer for a one-node tree (Example~3). \emph{Step}: assume
$\mathcal S(L)$ and $\mathcal S(R)$ are correct. $\mathit{lo}$ and
$\mathit{hi}$ extend by definition of min/max over a set (the set is the
union of $L$'s, $R$'s and $\{v\!\to\!\mathit{val}\}$); $\mathit{sz}$ is
correct by Ch.1's $\size(v)=1+\sum_c\size(c)$ induction;
$\mathit{isBST}(v)$ is correct because BST-ness at $v$ is the conjunction of
BST-ness on both sides and the two boundary checks $\mathit{hi}(L)<v$ and
$v<\mathit{lo}(R)$, and by induction the boundary values $\mathit{hi}(L)$,
$\mathit{lo}(R)$ are the actual maxima/minima of those subtrees.
$\mathit{best}(v)$ is the max of three candidates: the BST-at-$v$ size
(if $v$ is BST), $\mathit{best}(L)$ (max BST strictly inside $L$'s subtree
by induction), $\mathit{best}(R)$ (similar). Every BST subtree of $v$'s
subtree is either exactly $v$'s subtree or lives entirely inside $L$ or
$R$; the $\max$ over these three is the correct $\max$. \emph{Answer}: the
root's subtree is the whole tree, so $\mathit{best}(r)$ is the largest BST
subtree of $T$. \qed
\end{proofbox}

\solstep{Implementation}
\begin{lstlisting}[caption={LeetCode 333 -- five-coordinate tuple, one post-order fold, \coderef{c10-fold} instantiated}, label={lst:c10-subtree-soln}]
class Solution {
    struct S { bool isBST; int lo, hi, sz, best; };
    static S merge(TreeNode* u, const S& L, const S& R) {
        S s;
        s.lo  = min({L.lo, u->val, R.lo});
        s.hi  = max({L.hi, u->val, R.hi});
        s.sz  = 1 + L.sz + R.sz;
        s.isBST = L.isBST && R.isBST
                  && (L.sz == 0 || L.hi <  u->val)
                  && (R.sz == 0 || u->val <  R.lo);
        s.best = max({L.best, R.best, s.isBST ? s.sz : 0});
        return s;
    }
    static S dfs(TreeNode* u) {
        if (!u) return S{true, INT_MAX, INT_MIN, 0, 0};   // nil convention
        return merge(u, dfs(u->left), dfs(u->right));
    }
public:
    int largestBSTSubtree(TreeNode* root) { return dfs(root).best; }
};
\end{lstlisting}

\solstep{Trace}
\begin{tracebox}[Official Example~1: tree \texttt{[10,5,15,1,8,null,7]} $\Rightarrow$ 3]
Post-order on this binary tree: 1, 8, 5, 7, 15, 10.\\
Node 1 (leaf): isBST=T, lo=hi=1, sz=1, best=1.\\
Node 8 (leaf): isBST=T, lo=hi=8, sz=1, best=1.\\
Node 5 (left=1, right=8): $L{.}\mathit{hi}=1<5$ and $5<R{.}\mathit{lo}=8$;
both children BST $\Rightarrow$ isBST=T; lo$=\min(1,5,8)=1$,
hi$=\max(1,5,8)=8$, sz$=3$, best$=\max(1,1,3)=3$.\\
Node 7 (leaf): isBST=T, lo=hi=7, sz=1, best=1.\\
Node 15 (left=nil, right=7): nil passes trivially; boundary test
$15<R{.}\mathit{lo}=7$ \emph{fails} $\Rightarrow$ isBST=F; lo$=7$, hi$=15$,
sz$=2$, best$=\max(0,1,0)=1$ (subtree rooted at~15 is not BST, but its
right leaf~7 still contributes best$=1$ upward).\\
Node 10 (left=5 with best$=3$, right=15 with isBST=F):
$R{.}\mathit{isBST}=F$ blocks $u$ $\Rightarrow$ isBST=F; lo$=1$, hi$=15$,
sz$=6$, best$=\max(3,1,0)=3$.\\
Answer at root: best$=3$. \checkmark\ matches LC's official Example~1;
Examples~2 and~3 are strictness and single-node sanity checks and both
fold through the same rule.
\end{tracebox}

\solstep{Complexity}
\begin{complexitybox}
Time $\bigO{n}$ (each node is folded once; the fold itself is $\bigO{1}$ at
a binary tree, since $\deg\le 2$ after the $\nil$ convention). Memory
$\bigO{h}$ for the call stack; the tuple is five $\bigO{1}$ fields per
frame. On a chain at $n=10^{4}$, LC's stack handles it; if you port this to
general trees at $n\ge 2\cdot 10^{5}$ substitute \coderef{c09-iter}'s
two-phase form. Verdict: \ac, and this is the pattern the section is
\emph{about} -- tuple design over an unchanged DFS.
\end{complexitybox}

\solstep{Pitfalls and Variants}
\begin{warningbox}
Three ways to lose this problem: (i)~drop the strictness (use $\le$/$\ge$
instead of $<$/$>$) and duplicate values in the tree then make a
non-BST pass (Example~2's exact failure); (ii)~forget the fifth
coordinate and read the answer at $r$, which returns ``is the whole tree
BST'' rather than ``largest BST subtree anywhere''; (iii)~on a nil child,
use $\mathit{lo}=\mathit{hi}=0$ (or any finite value) and let a
real-valued node $<0$ fail the boundary test -- \emph{nil must be
$(\mathit{lo}{=}{+}\infty,\mathit{hi}{=}{-}\infty,\mathit{sz}{=}0)$}, so
the $\mathit{sz}{=}0$ short-circuits and the $\infty$ values are inert.
Variant worth doing: LC~2581 ``count ways where you can find the best
robbers'' (Q below) keeps the same ``extend tuple with a global-reading
coordinate'' idea but swaps BST-ness for ``maximum-weight independent
set on a tree'' -- that shape is exactly \secref{c10-twostate}'s, so it
pays to do it \emph{after} reading 13.2.
\end{warningbox}

\subsection{Practice Problems}

\begin{stuckbox}
\begin{itemize}[nosep, leftmargin=1.4em]
  \item Did your tuple carry a ``best'' coordinate so the answer is
        readable at the root? A global counter works but hides the DP
        structure.
  \item Does your nil-case initialise $\mathit{lo}=+\infty$,
        $\mathit{hi}=-\infty$, $\mathit{sz}=0$, $\mathit{isBST}{=}\text{true}$
        \emph{and} $\mathit{best}=0$? Every one matters; three of the four
        are silent failure modes.
  \item Are the boundary inequalities strict? Test on \texttt{[2,2]} (two
        equal nodes) -- the correct answer is $1$, and a
        non-strict version returns $2$.
\end{itemize}
\end{stuckbox}

\begin{practiceset}
  \practice{Count Ways Where You Can Find the Best Robbers}{LeetCode 2581}
    {https://leetcode.com/problems/count-ways-where-you-can-find-the-best-robbers/}{LC Hard L5, premium-era}
    {You \emph{own} this one next section. Read the statement, spot the
    tuple you would need, and stop. The answer will be far more legible after
    \secref{c10-twostate}'s in/out pair -- LC~2581 is a ``max weight
    independent set on a tree, plus a count of the optima'' fold, which is
    exactly \probref{c10-subtree-solved}'s tuple design pattern wearing the
    two-state shape. For now: attempt only the tuple \emph{design}, not the
    code.}
\end{practiceset}

\paragraph{Back-cites (URLs live elsewhere in the book).}
\begin{itemize}[nosep,leftmargin=1.4em]
  \item \textbf{LC~2049 ``count nodes with the highest score''} (\chref{c04}
        owns this URL): tuple is $(\mathit{sz},\mathit{left},\mathit{right},
        \mathit{score})$, fold is \coderef{c10-fold} plus Ch.4's parent-side
        trick that turns a rooted view into an unrooted one. Read
        \secref{c10-harvest} for why ``score'' is really a peak-decomposition.
  \item \textbf{LC~687 ``longest uni-value path''} (\chref{c02} owns this
        URL): Ch.2 solved it on \emph{binary} trees with a scalar fold
        carrying ``longest arm with same value starting at $v$''; the
        general-tree upgrade adds one coordinate ($\mathit{val}(v)$) so
        the fold can short-circuit on a value change. The pattern --
        ``carry the arm, close the path'' -- becomes \thmref{c09-dpjoin}
        in Ch.12 and \propref{c10-pairtop} in \secref{c10-harvest}.
\end{itemize}

\begin{recallbox}[10-second refresher -- 13.1 Subtree DP: state per node]
Each node carries one tuple $\mathcal S(v)$ computed by folding
$\{\mathcal S(c)\}$ (\defref{c10-subtree-state}). \lstinline{fold_dfs} is
\coderef{c09-dfs3} with a different payload; the induction is
\thmref{c01-height-recursion} and \thmref{c10-subtree-template} verbatim.
The design question is \emph{tuple size}: too small breaks closure, too big
breaks the budget. LC~333 (\probref{c10-subtree-solved}) shows the
five-coordinate tuple that makes BST-ness compose; Ch.12's
$\mathit{LD}$ table reappears here as a two-coordinate instance.
\selftest{
  \item State the tuple $\mathcal S(v)$ for ``subtree sum of node values'',
        and write its fold rule.
  \item Why is $\mathit{lo}(\nil)=+\infty$ and $\mathit{hi}(\nil)=-\infty$
        the right initialiser for LC~333, not $\mathit{lo}(\nil)=\mathit{hi}(\nil)=0$?
  \item Which coordinate makes the answer readable at $r$ for
        ``largest BST \emph{anywhere}'' (Trap~2)?
  \item Give one line for \thmref{c10-subtree-template}'s stack bound and
        one sentence for when it becomes an active problem.
  \item What breaks if a subtree's tuple fails to carry $\mathit{best}$ and
        you use a global instead? (Hint: rerooting later.)
}
\end{recallbox}
